
\SourceLeaf{397}%
\clearpage
\renewcommand{\BookRunningTitle}{Zur Theorie der complexen Zahlen}
\pdfbookmark[0]{Zur Theorie der complexen Zahlen}{nachlass-complexe-zahlen}
\begin{center}\Large ZUR THEORIE DER COMPLEXEN ZAHLEN.\end{center}
\BookRule{.15\textwidth}
\pdfbookmark[1]{I. Neue Theorie der Zerlegung der Cuben}{nachlass-complexe-I}
\begin{center}[I.]\par\medskip\large NEUE THEORIE DER ZERLEGUNG DER CUBEN.\end{center}

I. Wir nehmen an, es gebe eine Auflösung der Gleichung $x^3+y^3+z^3=0$, nemlich $x=a,\ y=b,\ z=c$, wo $a,b,c$ keinen gemeinschaftlichen Divisor haben, folglich auch unter sich Primzahlen sind. Wir setzen
\[
\begin{aligned}
b+c&=\alpha\\
c+a&=\beta\\
a+b&=\gamma
\end{aligned}
\]
wo nothwendig auch $\alpha,\beta,\gamma$ unter sich Primzahlen sein werden. Hätten nemlich $\alpha$ und $\beta$ einen gemeinschaftlichen Divisor, so würde dieser auch $a^3$ und $b^3$ messen, es müssten daher auch $a$ und $b$ einen gemeinschaftlichen Divisor haben.

Wir werden nun haben
\[
(\beta+\gamma-\alpha)^3+(\gamma+\alpha-\beta)^3+(\alpha+\beta-\gamma)^3=0
\]
allein es ist identisch
\[
\begin{aligned}
&(\beta+\gamma-\alpha)^3+(\gamma+\alpha-\beta)^3+(\alpha+\beta-\gamma)^3\\
&=(\alpha+\beta+\gamma)^3-24\alpha\beta\gamma
\end{aligned}
\]
Es wird folglich
\[
(\alpha+\beta+\gamma)^3=24\alpha\beta\gamma
\]

\SourceLeaf{398}%
Sind $\alpha,\beta,\gamma$ reelle Zahlen, so wird $\alpha+\beta+\gamma$ durch $3$ theilbar sein, also $\alpha+\beta+\gamma)^3$ durch $27$, folglich $\alpha\beta\gamma$ durch $9$. Es muss daher eine der Zahlen $\alpha,\beta,\gamma$ z.\ B.\ $\gamma$ durch $9$ theilbar sein, also $c^3$ ebenfalls, folglich $c$ durch $3$.

Sind hingegen $\alpha,\beta,\gamma$ imaginäre Zahlen, so schliessen wir, dass $\alpha+\beta+\gamma$ durch $1-\varepsilon$, folglich $24\alpha\beta\gamma$ durch $(1-\varepsilon)^3$, mithin $\alpha\beta\gamma$ durch $1-\varepsilon$ theilbar sein müsse. Es ist also eine der Zahlen $\alpha,\beta,\gamma$ durch $1-\varepsilon$ theilbar und folglich auch eine der Zahlen $a,b,c$.

II. Wir haben allgemein die identische Gleichung
\[
\begin{aligned}
&(p+q+r)^3+(p+q\varepsilon+r\varepsilon\varepsilon)^8+(p+q\varepsilon\varepsilon+r\varepsilon)^3\\
&=27pqr+3(p+q+r)(p+q\varepsilon+r\varepsilon\varepsilon)(p+q\varepsilon\varepsilon+r\varepsilon)
\end{aligned}
\]
Ist folglich $p+q+r=0$, so wird
\[
(p+q\varepsilon+r\varepsilon\varepsilon)^3+(p+q\varepsilon\varepsilon+r\varepsilon)^3-27pqr=0
\]

Sind hier $p,q,r,$ selbst Cuben, nemlich resp.\ $=a^3,b^3,c^3$; d.\ i.\ existirt eine Auflösung der Gleichung $x^3+y^3+z^3=0$, so wird
\[
\begin{aligned}
a^3+b^3\varepsilon+c^3\varepsilon\varepsilon&=a'\\
a^3+b^3\varepsilon\varepsilon+c^3\varepsilon&=b'\\
-3abc&=c'
\end{aligned}
\]
gesetzt, auch $a'^3+b'^3+c'^3=0$ werden. Aus dieser neuen Auflösung kann man auf gleiche Weise eine dritte ableiten u.\ s.\ w. Man überzeugt sich leicht, dass wenn die erste Auflösung in reellen Zahlen ist, auch die dritte eine solche sein wird.

Es ist noch zu bemerken, dass wenn $a,b,c$ keinen Factor gemein haben, dasselbe auch von $a',b',c'$ gelten wird, den Factor $1-\varepsilon$ abgerechnet. Es ist nemlich
\[
\begin{aligned}
\frac{a'}{1-\varepsilon}&=-\varepsilon\varepsilon a^3+\varepsilon b^3=a^2-\varepsilon c^3=-b^3+\varepsilon\varepsilon c^3\\
\frac{b'}{1-\varepsilon}&=a^3-\varepsilon b^3=-\varepsilon\varepsilon a^3+\varepsilon c^3=\varepsilon\varepsilon b^3-c^3\\
\frac{c'}{1-\varepsilon}&=(\varepsilon\varepsilon-1)abc
\end{aligned}
\]
Die beiden ersten Zahlen haben also weder mit $a$, noch mit $b$, noch mit $c$ einen Factor gemein, können auch nicht durch $1-\varepsilon$ theilbar sein, wenn nicht $a,b,c$
\SourceLeaf{399}%
zugleich durch $1-\varepsilon$ theilbar sind: daher haben jene auch keinen Factor mit der dritten gemein.

III. Aber auch der umgekehrte Weg wird offen stehen. Wir haben gesehen, dass eine der Grössen durch $1-\varepsilon$ theilbar ist: dies mag $c$ sein. Da man statt $a$ auch $a\varepsilon$ oder $a\varepsilon\varepsilon$ substituiren kann, und ebenso statt $b$ auch $b\varepsilon$ oder $b\varepsilon\varepsilon$, so dürfen wir voraussetzen, dass $a$ entweder $\equiv1$ oder $\equiv-1$ sein wird; wir werden das erstere voraussetzen, da im andern Fall $b\equiv1$ sein würde und nur mit $a$ vertauscht zu werden brauchte. Wir setzen demnach
\[
\begin{aligned}
a&=1+3\alpha\\
b&=-1+3\beta
\end{aligned}
\]
und
\[
\begin{aligned}
\frac{a\varepsilon+b\varepsilon\varepsilon}{\varepsilon-\varepsilon\varepsilon}&=1+(\varepsilon\varepsilon-\varepsilon)(\alpha\varepsilon+\beta\varepsilon\varepsilon)=A\\
\frac{a\varepsilon\varepsilon+b\varepsilon}{\varepsilon-\varepsilon\varepsilon}&=-1+(\varepsilon\varepsilon-\varepsilon)(\alpha\varepsilon\varepsilon+\beta\varepsilon)=B\\
\frac{a+b}{\varepsilon-\varepsilon\varepsilon}&=(\varepsilon\varepsilon-\varepsilon)(\alpha+\beta)=C
\end{aligned}
\]
wo $A+B+C=0$ wird, und $ABC=\frac{a^3+b^3}{(\varepsilon-\varepsilon\varepsilon)^3}=\left(\frac c{\varepsilon\varepsilon-\varepsilon}\right)^3$

Da hier
\[
\begin{aligned}
a&=-\varepsilon A+\varepsilon\varepsilon B\\
b&=\varepsilon\varepsilon A-\varepsilon B
\end{aligned}
\]
so können $A$ und $B$ keinen Factor gemein haben, weil ein solcher sonst auch gemeinschaftlicher Factor von $a$ und $b$ sein würde. Wegen $A+B+C=0$ kann folglich auch $C$ keinen Factor weder mit $A$ noch mit $B$ gemein haben. Hieraus folgt leicht, dass $A$ und $B$ und mithin auch $C$ Cuben sind. Denn $\left(\frac c{\varepsilon\varepsilon-\varepsilon}\right)^3$ wird durch $\varepsilon-\varepsilon\varepsilon$, folglich auch durch $(\varepsilon-\varepsilon\varepsilon)^3$ theilbar sein oder $\alpha+\beta$ durch $3$, daher wird $A\equiv1$, $B\equiv-1\ (\text{mod. }3)$.

Setzen wir nun
\[
\begin{aligned}
A&=a'^3\\
B&=b'^3\\
C&=c'^3
\end{aligned}
\]
so haben wir aus der Auflösung der Gleichung $x^3+y^3+z^3=0$

\SourceLeaf{400}%
\[
\begin{aligned}
x&=a\\
y&=b\\
z&=c
\end{aligned}
\]
eine andere abgeleitet
\[
\begin{aligned}
x&=a'\\
y&=b'\\
z&=c'
\end{aligned}
\]
\[
\text{wo }a'^3b'^3c'^3=\frac{c^3}{(\varepsilon\varepsilon-\varepsilon)^3}
\]
wo folglich $c'$ den Factor $1-\varepsilon$ einmal weniger enthalten wird, als $c$. Dies ist aber absurd, wenn $c$ nur durch eine bestimmte Potenz von $1-\varepsilon$ theilbar, d.\ i.\ wenn $c$ von $0$ verschieden ist. Denn durch Fortsetzung dieser Operationen würde man sonst am Ende auf eine Auflösung kommen, wo $z$ gar nicht durch $1-\varepsilon$ theilbar wäre gegen (I).
\BookRule{.3\textwidth}

Einen ähnlichen Weg kann man für die $5^{\mathrm{ten}}$ Potenzen nehmen. Ist nemlich $a^5+b^5+c^5=0$, so setzt man $b+c=\alpha,\ c+a=\beta,\ a+b=\gamma$, so wird
\[
\begin{aligned}
0&=(2a)^5+(2b)^5+(2c)^5\\
&=(\beta+\gamma-\alpha)^5+(\gamma+\alpha-\beta)^5+(\alpha+\beta-\gamma)^5\\
&=(\alpha+\beta+\gamma)^5-80\alpha\beta\gamma(\alpha\alpha+\beta\beta+\gamma\gamma)
\end{aligned}
\]

Es kann aber nicht $(\alpha+\beta+\gamma)^5=80\alpha\beta\gamma(\alpha\alpha+\beta\beta+\gamma\gamma)$ werden ohne dass eine der Zahlen $\alpha,\beta,\gamma$ durch $1-\varepsilon$ theilbar sei. Denn wären sie alle nicht theilbar, so müsste sowohl $\alpha+\beta+\gamma$ als $\alpha\alpha+\beta\beta+\gamma\gamma$ durch $1-\varepsilon$ theilbar sein, folglich auch $2(\alpha\alpha+\beta\beta+\gamma\gamma)+2(\alpha+\beta+\gamma)(\alpha+\beta-\gamma)=(2\alpha+\beta)^2+3\beta\beta$, was unmöglich ist.

Man kann dies auch so darstellen. Ist $a^5+b^5+c^5=0$, so wird
\[
\begin{aligned}
4(a+b+c)^5&=5(b+c)(c+a)(a+b)\\
&\qquad[(a+2b+3c)^2+3(a+c)^2-8(a+b+c)c]\\
&=5(b+c)(c+a)(a+b)\\
&\qquad[(b-c)^2+3(b+c)^2+4(a+b+c)a]
\end{aligned}
\]
\[
4(a+b+c)^5+5abc[(b-c)^2+3(b+c)^2]=5(a+b+c)\{\dotsfour\}
\]
Uebrigens würde der Beweis dem vorigen sehr ähnlich.

Versucht man aber denselben Gang bei den siebenten Potenzen, so gelingt es nicht zu beweisen, dass bei einer gegebenen Auflösung
\SourceLeaf{401}%
\[
a^7+b^7+c^7=0
\]
nothwendig eine der Grössen $a,b,c$ durch $7$ theilbar sein müsse. Es folgt nemlich nur
\[
\begin{aligned}
(\alpha+\beta+\gamma)^7={}&5\alpha\beta\gamma\{3(\alpha^4+\beta^4+\gamma^4)\\
&\qquad+10(\alpha\alpha\beta\beta+\alpha\alpha\gamma\gamma+\beta\beta\gamma\gamma)\}
\end{aligned}
\]
welches bestehen kann, ohne dass $\alpha,\beta,\gamma$ durch $1-\varepsilon$ theilbar wäre.

Hoffentlich wird sich indessen dies in Zukunft aus der Natur der Determinanten und der Einheitszahlen ableiten lassen.
\BookRule{.3\textwidth}

